% Generated draft: requires mathematical and source review.
\documentclass[11pt]{article}
\usepackage[margin=25mm]{geometry}
\usepackage{fontspec}
\setmainfont{Latin Modern Roman}
\usepackage{amsmath,amssymb,mathtools,unicode-math}
\setmathfont{Latin Modern Math}
\usepackage{enumitem,needspace,xurl,hyperref}
\setlength{\parindent}{0pt}
\setlength{\parskip}{5pt}
\setlength{\emergencystretch}{3em}
\newcommand{\N}{\mathbb{N}}
\newcommand{\Z}{\mathbb{Z}}
\newcommand{\Q}{\mathbb{Q}}
\newcommand{\R}{\mathbb{R}}
\newcommand{\C}{\mathbb{C}}
\newcommand{\F}{\mathbb{F}}
\newcommand{\A}{\mathbb{A}}
\newcommand{\PP}{\mathbb{P}}
\newcommand{\E}{\mathbb{E}}
\newcommand{\eps}{\varepsilon}
\newcommand{\dd}{\,\mathrm{d}}
\begin{document}
\section*{UnsolvedMath ID: 1872}
\textbf{Source problem code:} catalogue entry
\subsection{Definitions and mathematical statement}
A prime is a positive integer $p\ge2$ whose only positive divisors are $1$ and $p$. For an integer $n\ge3$, consider the property
\[
 \forall k\in\Z_{\ge1},\qquad 2^k<n\ \Longrightarrow\ n-2^k\text{ is prime}.
\]
The nonvacuous finite-list conjecture attributed in the source to Erd\H{o}s is
\[
 \bigl\{n\in\Z_{\ge3}:\ \forall k\ge1,\
 2^k<n\Longrightarrow n-2^k\text{ is prime}\bigr\}
 =\{4,7,15,21,45,75,105\}.
\]
The restriction $k\ge1$ implements the source's $2\le2^k<n$, so the power $2^0=1$ is excluded. The lower bound $n\ge3$ is an explicitly identified nonvacuous interpretation: the source does not state a domain or lower bound on $n$. If all positive integers are allowed literally, $n=1,2$ satisfy the universal condition vacuously and are additional trivial values. If arbitrary integers are allowed, there are still further vacuous values. Thus the exact finite list requires a source-confirmed nonvacuity convention rather than an implicit omission of these cases.

The background also records a separate Erd\H{o}s conjecture. A positive integer $m$ is squarefree if $p^2\nmid m$ for every prime $p$. This companion assertion is that there are infinitely many positive integers $n$ such that
\[
 \forall k\in\Z_{\ge0},\qquad 2^k<n\ \Longrightarrow\
 n-2^k\text{ is squarefree}.
\]
Here $k=0$ is included, because the background uses $1\le2^k<n$. Equivalently, for every bound $N$ there should be an $n>N$ satisfying these squarefreeness conditions. It is not the same assertion as completeness of the prime-value list.

The supplied literature triage checked 17 August 2026 reports that a modern compendium still lists the prime-value question as open, and that no proof of list completeness or additional nontrivial example was verified. It supplies no distinct mathematical partial theorem and no separate updated assessment of the companion squarefree conjecture. The registry's main label and triage label are open; these are dated source assessments.
\subsection{Short English statement}
Are $4,7,15,21,45,75,105$ the complete nontrivial list of integers whose differences from every smaller power $2^k\ge2$ are prime? The source must exclude vacuous tiny values. Its background separately asks for infinitely many integers whose differences from every smaller power of two, including $1$, are squarefree.
\subsection{Sources}
\begin{itemize}
\item Local UnsolvedMath registry, permanent ID $1872$, source code \texttt{catalogue entry}, ``Erd\H{o}s Conjecture on $n-2^k$ Prime,'' including the background and literature triage checked 17 August 2026.
\item Richard Guy, \emph{Unsolved Problems in Number Theory}, Section A19, as identified in the supplied background; no edition or page is supplied. Both questions are attributed there to Erd\H{o}s.
\item P. Moree, \emph{Artin's primitive root conjecture---a survey}, MPIM preprint 2012-53, \url{https://archive.mpim-bonn.mpg.de/1236/1/preprint_2012_53.pdf}; the supplied triage cites it as retaining the Erd\H{o}s prime-value question.
\end{itemize}
\end{document}
